\begin{examplebox}
\textbf{例1｜最快识别：只求内心横坐标}

已知双曲线 $\frac{x^2}{9}-\frac{y^2}{16}=1$，焦点为 $F_1,F_2$。
当 $P$ 在右支或左支时，$\triangle PF_1F_2$ 的内心分别落在哪条直线上？

\textbf{解析：} $a=3$，由 $PF_1-PF_2=\pm2a$ 及
$x_I=(PF_1-PF_2)/2$，得
\[
\boxed{\text{右支 }x_I=3;\qquad \text{左支 }x_I=-3.}
\]
\textbf{识别信号：} “焦点三角形内心横坐标”“定直线”“顶点切线”，立即想到\textbf{焦半径差}，不必设内心求方程。

\tcbline
\textbf{例2｜四心全求，并用欧拉线反向验算}

双曲线 $\frac{x^2}{9}-\frac{y^2}{16}=1$ 上取
$P\left(5,\frac{16}{3}\right)$，求 $G,I,H,O'$。

\textbf{解析：} $a=3,b=4,c=5,e=5/3$。
先检验 $25/9-(256/9)/16=1$，因此 $P$ 在右支上。
\begin{align*}
G&=\left(\frac53,\frac{16}{9}\right),\\
I&=\left(3,\frac{3\cdot(16/3)}{3+5}\right)=(3,2),\\
H&=\left(5,\frac{25-25}{16/3}\right)=(5,0),\\
O'&=\left(0,\frac{25+256/9-25}{32/3}\right)
=\left(0,\frac83\right).
\end{align*}
检验欧拉线：
\[
3G-2O'=(5,16/3)-(0,16/3)=(5,0)=H.
\]
\textbf{几何解释：} $x_0=c=5$，直线 $PF_2$ 竖直、底边 $F_1F_2$ 水平，
三角形在 $F_2$ 处为直角，因此垂心恰好就是 $F_2$。

\textbf{变式：} 保持纵坐标，把 $P$ 换成 $\left(-5,\frac{16}{3}\right)$。
由关于 $y$ 轴的对称性，内心立即是 $\boxed{(-3,2)}$，
垂心为 $(-5,0)$，外心仍为 $(0,8/3)$。

\tcbline
\textbf{例3｜最常见陷阱：面积比的正确结果}

双曲线右支上 $P(x_0,y_0)$ 构成焦点三角形，$I$ 为内心。
求 $S_{\triangle IF_1F_2}/S_{\triangle PF_1F_2}$。

\textbf{解析：} 两个三角形的底边相同，面积比等于对底边的高之比：
\[
\frac{S_{\triangle IF_1F_2}}{S_{\triangle PF_1F_2}}
=\frac{|y_I|}{|y_0|}
=\frac{c}{c+ex_0}
=\boxed{\frac{a}{a+x_0}}.
\]
\textbf{注意：} 最后一步用 $e=c/a$，所以
$c+ex_0=c(a+x_0)/a$，\textbf{并不是} $a+ex_0$。
若不限定分支，正确统一式为 $\boxed{a/(a+|x_0|)}$。

代入例2的点检验：$|y_I|/|y_0|=2/(16/3)=3/8$，
且 $a/(a+x_0)=3/(3+5)=3/8$。

\tcbline
\textbf{例4｜完整轨迹：为什么不是两条无限长直线？}

求双曲线 $\frac{x^2}{a^2}-\frac{y^2}{b^2}=1$ 上
非顶点动点 $P$ 的焦点三角形内心轨迹。

\textbf{解析：} 已知 $x_I=\pm a$。但还要确定纵坐标的范围。
设 $u=|x_0|>a$，由双曲线方程
$|y_0|=\frac ba\sqrt{u^2-a^2}$，于是
\[
|y_I|=\frac{a|y_0|}{a+u}
=b\sqrt{\frac{u-a}{u+a}}.
\]
因为 $u>a$，所以 $0<|y_I|<b$。
而 $u$ 从 $a$ 增大至无穷时，右边连续增大并遍历 $(0,b)$。
故轨迹准确写成
\[
\boxed{x=\pm a,\qquad 0<|y|<b.}
\]
\textbf{易错：} 不能把 $x=\pm a$ 全部看成真实轨迹；
不能包含 $(\pm a,0)$ 或 $(\pm a,\pm b)$。

\tcbline
\textbf{例5｜逆向思考：已知内心，反求动点}

在 $\frac{x^2}{9}-\frac{y^2}{16}=1$ 上，
若 $\triangle PF_1F_2$ 的内心是 $I(3,2)$，求 $P$。

\textbf{解析：} $I_x=3=a$ 表明 $P$ 在右支，令 $x_0=u>3$。
由例4的表达式，
\[
2=4\sqrt{\frac{u-3}{u+3}}
\Longrightarrow\frac14=\frac{u-3}{u+3}
\Longrightarrow u=5.
\]
再由 $y_I=3y_0/(3+u)$，得到
\[
y_0=\frac{2(3+5)}3=\frac{16}{3},\qquad
\boxed{P\left(5,\frac{16}{3}\right)}.
\]
\textbf{推广：} 若 $I=(\varepsilon a,t)$、$0<|t|<b$，则
\[
\boxed{x_0=\varepsilon a\frac{b^2+t^2}{b^2-t^2}},\qquad
\boxed{y_0=\frac{2b^2t}{b^2-t^2}}.
\]
这是把轨迹结论转化成求动点坐标的反向公式。

\tcbline
\textbf{例6｜高价值迁移：椭圆与双曲线为什么结果不同？}

设焦点仍为 $F_1(-c,0),F_2(c,0)$，比较下列两种情况：
\begin{enumerate}
  \item 椭圆 $x^2/a^2+y^2/b^2=1$，$a>b>0$、$c^2=a^2-b^2$；
  \item 双曲线 $x^2/a^2-y^2/b^2=1$，$c^2=a^2+b^2$。
\end{enumerate}

\textbf{解析：} 对任何不退化的焦点三角形，
\[
x_I=\frac{PF_1-PF_2}{2},\qquad
 y_I=\frac{cy_0}{c+(PF_1+PF_2)/2}.
\]
椭圆中 $PF_1-PF_2=2ex_0$、$PF_1+PF_2=2a$，所以
\[
\boxed{I_{\text{椭圆}}\left(ex_0,\frac{cy_0}{a+c}\right)}.
\]
双曲线中 $PF_1-PF_2=\pm2a$、$PF_1+PF_2=2e|x_0|$，所以
\[
\boxed{I_{\text{双曲线}}\left(\pm a,\frac{cy_0}{c+e|x_0|}\right)}.
\]
\textbf{迁移要点：} 同一几何母公式，先判断固定的是焦半径\textbf{和}还是\textbf{差}，
就能快速得到不同的内心性质。
\end{examplebox}
